\section{Monetary policy and the adoption margin}
\label{sec:monetary}

Two questions connect monetary policy to the adoption margin, and they point in
opposite directions. Does the margin reshape how a monetary shock transmits? And
does the monetary stance, during an energy shock, reshape the adoption wave? The
first is a robustness check; the second is where the interest-rate sensitivity of
green durables that \citet{Dietrich2025} identify in a representative agent shows up
in our setting.

\subsection{The adoption margin does not reshape monetary transmission}
\label{sec:mp_transmission}

The energy-shock experiments hold monetary policy at the AMRS constant-real-rate
baseline. A $100$-basis-point (annualised) tightening isolates the role of the
rule itself. Figure~\ref{fig:monetary} contrasts the constant-real-rate baseline
with a standard Taylor rule $(\phi_\pi,\phi_{\pi e})=(1.5,0)$. The response is the
textbook one, output and consumption fall, inflation falls, and it is
\emph{not} materially reshaped by the adoption margin: at a $5\%$ green steady
state the durable-choice channel is too small to move aggregate monetary
transmission (on impact, output $-0.85$ under the Taylor rule vs.\ $-1.13$ under
the constant-real-rate rule; inflation $-0.43$ vs.\ $-0.50$ annualised). The
adoption margin matters for how the economy absorbs an \emph{energy} shock, not
for how it transmits a \emph{monetary} one.

\begin{figure}[H]\centering
\IfFileExists{output/fig_monetary_import.pdf}{\includegraphics[width=\textwidth]{output/fig_monetary_import.pdf}}{\textit{[figure output/fig\_monetary\_import.pdf: run run\_summary\_table.py]}}
\caption{Monetary-policy shock ($100$bp annualised tightening): Taylor rule
(solid) versus the constant-real-rate baseline (dashed). Nominal rate, output,
inflation, consumption.}
\label{fig:monetary}
\end{figure}

\subsection{The monetary stance moves the wave, but weakly}
\label{sec:mp_stance}

The reverse question is Dietrich's. In their representative-agent economy green
durables are interest-rate sensitive, so a strict inflation target slows the
transition. The mechanism survives here and acquires a distributional edge: the
switch is a lumpy purchase financed against the borrowing constraint, so its user
cost rises with the real rate precisely for the constrained households that account
for most of the crisis adoption (Figure~\ref{fig:switch_heatmap}).

The stance operates on the wave, but weakly. Table~\ref{tab:monetary_stance} holds
the brown-price shock fixed and varies the rule. A Taylor rule that leans against
the energy-driven inflation ($\phi_\pi=1.5$) raises the ex-ante real rate, $+5.3$
cumulated over $H=24$,and shrinks the peak green-share response from $8.9$ to
$8.6$ pp. A deliberate $25$-basis-point accommodation does the opposite, lifting it
to $8.9$ pp and cutting the output loss; the standalone rate cut raises adoption on
its own ($D^G$ peak $+0.09$ pp, consumption $+5.9$ cumulated) as easier financing
relaxes the constraint on the switch. The sign is Dietrich's; the magnitude is
small.

\input{output/tab_monetary_stance_import.tex}

Figure~\ref{fig:monetary_stance} sweeps the inflation coefficient. Stricter
targeting shrinks the wave monotonically, from $8.7$ pp at $\phi_\pi=1.25$ to $8.4$
at $\phi_\pi=3$,but the whole range spans half a percentage point, while the
cumulative output loss grows from $-77$ to $-103$. Strict inflation targeting during
an energy crisis buys a marginal slowdown of the transition at a first-order output
cost.

\begin{figure}[H]\centering
\IfFileExists{output/fig_monetary_stance_import.pdf}{\includegraphics[width=\textwidth]{output/fig_monetary_stance_import.pdf}}{\textit{[figure output/fig\_monetary\_stance\_import.pdf: run run\_monetary\_stance.py]}}
\caption{The monetary stance and the adoption wave under the brown-price shock, as
the Taylor inflation coefficient $\phi_\pi$ rises (with $\phi_{\pi e}=0$). Left:
peak $\Delta D^G$. Right: cumulative output. The dotted line is the
constant-real-rate baseline. Stricter targeting shrinks the wave marginally and
deepens the output loss substantially.}
\label{fig:monetary_stance}
\end{figure}

The comparison that matters is with the fiscal instruments. The peak wave's
semi-elasticity to the real rate is about $-0.05$ pp per unit of cumulative rate, so
moving the wave by one percentage point would take a cumulative rate change of order
twenty, an implausible monetary action, whereas the retail price cap moves it by
$8.9$ pp at a stroke (Section~\ref{sec:cap}). Monetary policy reaches the adoption
margin only through the financing cost of the switch, a second-order channel; the
cap and the transfer reach it through the relative price itself, which is
first-order. This is why we keep monetary policy at the constant-real-rate baseline
in the main text, not because the stance is neutral for adoption, but because it is
dominated, by an order of magnitude, by the instruments the paper ranks.

\paragraph{An accommodative offset.} A natural crisis exercise runs an accommodative
innovation \emph{alongside} the energy shock, leaning against it rather than waiting.
Figure~\ref{fig:monetary_offset} sizes the accommodation to halve the impact output
loss (about a $40$-basis-point cut). It cushions output (cumulative $-61$ to $-43$)
and consumption ($-106$ to $-88$) and lifts the adoption wave ($8.9$ to $9.2$ pp),
the financing channel working in its favour, but pays for the relief with more
inflation (impact annualised inflation $+7.3\%$ to $+9.0\%$). The offset buys output
stabilisation and a small transition dividend at the cost of the price level; it does
not substitute for the relative-price instruments, which stabilise output without
that inflation cost.

\begin{figure}[H]\centering
\IfFileExists{output/fig_monetary_offset_import.pdf}{\includegraphics[width=\textwidth]{output/fig_monetary_offset_import.pdf}}{\textit{[figure output/fig\_monetary\_offset\_import.pdf: run run\_monetary\_offset.py]}}
\caption{Accommodative monetary offset run alongside the brown-price shock, sized to
halve the impact output loss: energy shock alone (solid) versus energy shock plus
accommodation (dashed). The offset cushions output and consumption and amplifies
adoption, at the cost of higher inflation. Level deviations $\times100$ (real rate in
levels).}
\label{fig:monetary_offset}
\end{figure}
